\documentclass[11pt]{article}
\usepackage[a4paper,margin=2.4cm]{geometry}
\usepackage{amsmath}
\usepackage{booktabs}
\usepackage{xcolor}
\usepackage[hidelinks]{hyperref}

\definecolor{accent}{HTML}{1F5A94}
\setlength{\parindent}{0pt}
\setlength{\parskip}{0.65em}

\begin{document}
\title{\textcolor{accent}{Tool Report}}
\author{Name -- Student ID}
\date{August 24, 2026}
\maketitle

\section{Result}
For a data set containing $n$ latency measurements $t_i$, the arithmetic
mean is defined by Equation~\eqref{eq:mean-latency}.
\begin{equation}
  \bar{t}=\frac{1}{n}\sum_{i=1}^{n} t_i
  \label{eq:mean-latency}
\end{equation}
The compact build check in Table~\ref{tab:build-check} records the two
essential parts of this one-page technical note.

\begin{table}[htbp]
  \centering
  \caption{Build check}
  \label{tab:build-check}
  \begin{tabular}{cc}
    \toprule
    Item & Status \\
    Formula reference & Passed \\
    \bottomrule
  \end{tabular}
\end{table}

The numbered formula and the captioned table are both referenced in the
body text. The document is therefore ready for a command-line build.
\end{document}
